\subsection{下降：平展空间的情形}

设 X 和 Y 是拓扑空间，设 \(f\colon X\to Y\) 是连续映射。设 T 和 T' 是 Y-空间；给 X-空间 \(X\times_{Y}T\) 和 \(X\times_{Y}T'\) 赋予其规范下降数据，并记 \(\mathcal{C}_{f}(X\times_{Y}T;X\times_{Y}T')\) 为从 \(X\times_{Y}T\) 到 \(X\times_{Y}T'\) 且与这些下降数据相容的 X-态射集合。对于每个 Y-态射 \(\varphi\colon T\to T'\)，X-态射 \(f^{*}(\varphi)\colon(x,t)\mapsto(x,\varphi(t))\) 与规范下降数据相容，并且它从 \(X\times_{Y}T\) 到 \(X\times_{Y}T'\)。将由此定义的映射记为 \(f^{*}\colon\mathcal{C}_{Y}(T;T')\to\mathcal{C}_{f}(X\times_{Y}T;X\times_{Y}T')\)。

命题 4. --- 假设映射 f 严格且满射，并且 T 是平展 Y-空间。那么，映射 \(f^{*}\)：\(\mathcal{C}_{\mathrm{Y}}(\mathrm{T};\mathrm{T}^{\prime}) \to \mathcal{C}_{f}(\mathrm{X} \times_{\mathrm{Y}} \mathrm{T}; \mathrm{X} \times_{\mathrm{Y}} \mathrm{T}^{\prime})\) 是双射。

记 \(\tau\)（分别记 \(\tau'\)）为与规范下降数据关联的 \(X\times_{Y}T\)（分别为 \(X\times_{Y}T'\)）上的等价关系。由于映射 f 是满射，投影 \(\mathrm{pr}_2\colon \mathrm{X}\times_{\mathrm{Y}}\mathrm{T} \to \mathrm{T}\) 是满射，且规范映射 \((\mathrm{X}\times_{\mathrm{Y}}\mathrm{T})/R_{\tau} \to \mathrm{T}\) 是双射。特别地，映射 \(f^*\) 是单射。现证其满射。设 \(\varphi\colon \mathrm{X}\times_{\mathrm{Y}}\mathrm{T} \to \mathrm{X}\times_{\mathrm{Y}}\mathrm{T}'\) 是与规范下降数据相容的 X-态射。因此，对于 \((x,t) \in \mathrm{X}\times_{\mathrm{Y}}\mathrm{T}\)，有 \(\varphi(x,t) = (x,\overline{\varphi}(t))\)，其中 \(\overline{\varphi}\) 是从 T 到 T' 的映射。

按 \(\overline{\varphi}\) 的定义，映射 \(\overline{\varphi} \circ \mathrm{pr}_2: \mathrm{X} \times_{\mathrm{Y}} \mathrm{T} \to \mathrm{T}'\) 等于 \(\mathrm{pr}_2 \circ \varphi\)，因而连续。由于映射 f 满射且严格，而 T 是 Y-平展空间，投影 \(\mathrm{pr}_2: \mathrm{X} \times_{\mathrm{Y}} \mathrm{T} \to \mathrm{T}\) 是严格的（I, 第~32~页，注 3）。根据 I, 第~18~页命题 9，映射 \(\overline{\varphi}\) 因而连续。它是满足 \(f^*(\overline{\varphi}) = \varphi\) 的 Y-态射，这就证明了命题。

推论。--- 设 X 和 Y 是拓扑空间，设 \(f: X \rightarrow Y\) 是严格且满射的映射。设 T 和 T' 是 Y-平展空间。若存在与规范下降数据相容的 X-同构 \(X \times_{Y} T \rightarrow X \times_{Y} T'\)，则 Y-空间 T 和 T' 同构。

设 \(\psi\colon X\times_{Y}T\to X\times_{Y}T'\) 是平展空间之间的 X-同构。根据命题 4，存在唯一的 Y-空间态射 \(\varphi\colon T \to T'\)，使得 \(\psi = f^{*}(\varphi)\)。由于 f 满射，映射 \(\varphi\) 是双射。由于 Y-空间 T 和 T' 都是平展的，根据 I, 第~30~页推论 2，\(\varphi\) 是同构。

命题 5. --- 设 X 和 Y 是拓扑空间，且 \(f: X \to Y\) 是连续映射。假设映射 f 是逆紧且分离的，或者是开映射。关于 f 在 X-平展空间上的每个下降数据都是有效的。此外，若 f 是满射，则商空间是 Y-平展空间。

设 \((Z,p)\) 是 X-平展空间，\(\tau\) 是关于 f 在 \((Z,p)\) 上的下降数据。记 \(R_{\tau}\) 为与 \(\tau\) 关联的 Z 上的等价关系，令 \(g: Z \to Z/R_{\tau}\) 为规范满射，并记 h 为从 Z 到 \(X \times_{Y} (Z/R_{\tau})\) 的由 \(z \mapsto (p(z), g(z))\) 定义的映射。它连续且双射；现证明它是同胚。

\begin{enumerate}
\def\labelenumi{\alph{enumi})}
\tightlist
\item
  首先假设映射 f 是开映射。
\end{enumerate}

设 \(z_0\) 是 Z 中一点；令 \(x_0 = p(z_0)\)。由于 p 是平展的，存在 X 中 \(x_0\) 的一个邻域 U，以及 p 在 U 上的连续截面 \(s: \mathrm{U} \to \mathrm{Z}\)，使得 \(s(x_0) = z_0\)。映射 \(p \times p: \mathrm{Z} \times_{\mathrm{Y}} \mathrm{Z} \to \mathrm{X} \times_{\mathrm{Y}} \mathrm{X}\) 是平展的，并且从 \(\mathrm{U} \times_{\mathrm{Y}} \mathrm{U}\) 到 \(\mathrm{Z} \times_{\mathrm{Y}} \mathrm{Z}\) 的两个映射分别由 \((x, x') \mapsto (s(x), s(x'))\) 和 \((x, x') \mapsto (s(x), \tau(s(x), x'))\) 定义，它们是 \(\mathrm{U} \times_{\mathrm{Y}} \mathrm{U}\) 上方的连续截面。由于它们在 \(\mathrm{U} \times_{\mathrm{Y}} \mathrm{U}\) 中形如 \((x, x)\) 的每一点处重合，所以它们在某个开邻域 V 上重合，该邻域包含 \(\Delta_{\mathrm{U}}\) 且包含于 \(\mathrm{U} \times_{\mathrm{Y}} \mathrm{U}\)。令 \(\mathrm{U}_0\) 是 X 中 \(x_0\) 的一个邻域，使得 \(\mathrm{U}_0 \times_{\mathrm{Y}} \mathrm{U}_0\) 包含于 V。设 \((x, u)\) 是 \(\mathrm{U}_0 \times_{\mathrm{Y}} g(s(\mathrm{U}_0))\) 中的一点；令 \(x' \in \mathrm{U}_0\)，使得 \(u = g(s(x'))\)。因此有 \(s(x) = \tau(s(x'), x)\)，从而 \(\mathrm{R}_{\tau}\{s(x), s(x')\}\)，且 \((x, u) = (x, g(s(x'))) = (x, g(s(x))) = h(s(x))\)。

由于映射 g 是开映射（IV, 第~386~页，命题 2），集合 \(U_{0} \times_{Y} g(s(U_{0}))\) 是 \(X \times_{Y}(Z/R_{\tau})\) 的开子集，在其上映射 \(h^{-1}\) 等于连续映射 \(s \circ pr_{1}\)。因此映射 h 是同胚。

\begin{enumerate}
\def\labelenumi{\alph{enumi})}
\setcounter{enumi}{1}
\tightlist
\item
  现在假设映射 f 逆紧且分离。设 \(z_{0}\) 是 Z 中一点；令 \(x_{0}=p(z_{0})\) 且 \(y_{0}=f(x_{0})\)。由 \(x\mapsto\tau(z_{0},x)\) 给出的映射 s 是 p 在 \(f^{-1}(y_{0})\) 上的截面。集合 \(f^{-1}(y_{0})\) 是紧的（TG, I, 第~75~页，定理 1 以及 I, 第~26~页，注 2），并且由于 f 分离，\(f^{-1}(y_{0})\) 中任意两个不同的点在 X 中有不交邻域（I, 第~25~页，命题 1）。根据 I, 第~37~页定理 2，存在 X 中的一个邻域 \(U_{0}\)，它包含 \(f^{-1}(y_{0})\)，以及截面 \(s_{0}\)，它是 p 在 \(U_{0}\) 上的截面并延拓 s。集合 \(s_{0}(U_{0})\) 是 Z 的开子集，因为 p 是平展的（I, 第~30~页，推论 3），并且包含集合 \(\{z_{0}\}\) 关于关系 \(R_{\tau}\) 的饱和。映射 g 是闭映射（IV, 第~386~页，命题 2）。因此存在 \(Z/R_{\tau}\) 的一个开集 V，使得 \(W = g^{-1}(V) \cap p^{-1}(U_{0})\) 是 \(z_{0}\) 的邻域且包含于 \(s_{0}(U_{0})\)（I, 第~75~页，引理）。
\end{enumerate}

现在设 \((x,u)\in\mathrm{U}_{0}\times_{\mathrm{Y}}\mathrm{V}\)，并设 z 是 Z 中满足 \(h(z)=(x,u)\) 的唯一点；根据定义，有 \(z\in W\)。由于 \(\mathrm{W}\subset s_{0}(\mathrm{U}_{0})\)，有 \(z=s_{0}(p(x))\)。这说明，\(h^{-1}\) 在开子集 \(U_{0}\times_{Y}V\)（它是 \(X\times_{Y}(Z/R_{\tau})\) 的开子集）上的限制等于 \(s_{0}\circ p\circ pr_{1}\)。因此，h 是同胚。

这样就证明下降数据 \(\tau\) 是有效的。映射 f 是万有严格的（I, 第~20~页，推论）。若 f 满射，则根据 I, 第~31~页命题 8，\(q: Z/R_{\tau} \to Y\) 是平展的。

